\chapter*{Introduction générale}
\addcontentsline{toc}{chapter}{Introduction générale}

Le secteur de la location de véhicules au Maroc, particulièrement dynamisé par la relance du tourisme international et l'expansion des zones d'affaires, connaît une mutation sans précédent. Face à l'émergence des plateformes mondiales de réservation (Booking, Expedia), les agences locales se doivent de digitaliser leurs processus pour survivre et rester compétitives. Cependant, une proportion alarmante d'agences continue de s'appuyer sur des méthodologies traditionnelles. C'est le cas de \textbf{Prestige Maroc}, une agence de location de renom, qui se trouvait freinée dans son expansion par une gestion logistique et administrative obsolète, reposant quasi-exclusivement sur des carnets de réservation manuscrits et des fichiers tableurs (Excel) isolés.

Le présent projet de fin d'études s'inscrit dans le cadre de l'obtention du diplôme de \textbf{Brevet de Technicien Supérieur (BTS) en Développement d'Applications Informatiques (DAI)} au sein de l'Établissement Alkendi. L'objectif fondamental qui m'a été confié était de concevoir, architecturer et déployer un \textbf{Écosystème Numérique Unifié}. Plus qu'un simple site web "vitrine", ce projet vise à remplacer la totalité du système d'information de l'entreprise par une plateforme web transactionnelle B2C complète, couplée à un outil logiciel de gestion d'entreprise (ERP/CRM) pour l'administration.

Afin de garantir la pérennité, la robustesse et l'évolutivité de ce système face à une potentielle charge de milliers d'utilisateurs simultanés en haute saison, nous avons écarté les architectures monolithiques traditionnelles. J'ai pris la décision stratégique d'implémenter une \textbf{Architecture Microservices} moderne, propulsée par la pile technologique \textbf{MERN} (MongoDB, Express, React, Node.js). 

Pour mener à bien cette ingénierie complexe dans les délais impartis, le projet a été piloté selon la méthodologie \textbf{Agile Scrum}, garantissant une livraison continue de fonctionnalités testées.

Ce rapport détaillera avec précision l'ensemble du cycle de vie du développement logiciel (SDLC). Nous débuterons par la présentation du contexte d'accueil et de l'entreprise (Chapitre 1), suivie de l'étude critique de l'existant (Chapitre 2). Nous dresserons ensuite un cahier des charges rigoureux (Chapitre 3) avant d'aborder la modélisation UML et l'analyse fonctionnelle (Chapitre 4). Les justifications de nos choix technologiques (Chapitre 5) introduiront la partie cœur de ce rapport : la réalisation concrète des microservices (Chapitre 6). Enfin, nous validerons l'intégrité de notre système par une campagne de tests exhaustifs et l'implémentation de pipelines d'intégration continue (Chapitre 7).
